\documentclass{article}
\usepackage{amsmath,amssymb,amsthm,mathtools}
\newtheorem{theorem}{Theorem}
\newtheorem{lemma}{Lemma}
\newtheorem{proposition}{Proposition}
\newtheorem{corollary}{Corollary}
\newtheorem{definition}{Definition}
\newtheorem{example}{Example}
\title{ShadowBench algebra/L2/alg\_gen\_L2\_006}
\begin{document}
\maketitle
% Problem id: algebra/L2/alg_gen_L2_006
% Companion instructions: docs/instructions.md
% Lean target: ShadowBench/Source/Main.lean

\begin{theorem}[ABM_algebra_L2_alg_gen_L2_006_item_1]
Prove that ideal-variety correspondence is inclusion-reversing, i.e., 
if $I_1 \subseteq I_2$ are ideals, then $\mathbf V(I_1) \supseteq \mathbf V(I_2)$, 
and similarly, if $V_1 \subseteq V_2$ are affine algebraic sets, then $\mathbf I(V_1) \supseteq \mathbf I(V_2)$. 
Also prove that $\mathbf V(\sqrt{I})=\mathbf V(I)$ for any ideal $I$.
\end{theorem}

\begin{proof}
First, let $I_1 \subseteq I_2$ be ideals. If $a \in \mathbf V(I_2)$, then $f(a)=0$ for every $f \in I_2$.
Since $I_1 \subseteq I_2$, we also have $f(a)=0$ for every $f \in I_1$, hence $a \in \mathbf V(I_1)$.
Therefore $\mathbf V(I_2) \subseteq \mathbf V(I_1)$, i.e., $\mathbf V(I_1) \supseteq \mathbf V(I_2)$.

Similarly, if $V_1 \subseteq V_2$ and $f \in \mathbf I(V_2)$, then $f(a)=0$ for every $a \in V_2$.
In particular, $f(a)=0$ for every $a \in V_1$, so $f \in \mathbf I(V_1)$.
Thus $\mathbf I(V_2) \subseteq \mathbf I(V_1)$, i.e., $\mathbf I(V_1) \supseteq \mathbf I(V_2)$.

Finally, we prove $\mathbf V(\sqrt{I})=\mathbf V(I)$.
Since $I \subseteq \sqrt{I}$, the inclusion-reversing property implies $\mathbf V(\sqrt{I}) \subseteq \mathbf V(I)$.
Conversely, let $a \in \mathbf V(I)$ and let $f \in \sqrt{I}$. Then $f^m \in I$ for some $m \ge 1$,
so $(f(a))^m = (f^m)(a) = 0$. Since we are over a field, this implies $f(a)=0$, hence $a \in \mathbf V(\sqrt{I})$.
Therefore $\mathbf V(I) \subseteq \mathbf V(\sqrt{I})$, and the equality follows.
\end{proof}
\end{document}
